\documentclass[twoside,11pt]{article}

\usepackage{enumitem}
\usepackage{graphicx,color}
\usepackage{hyperref}
\usepackage{url}
\usepackage{booktabs}
\usepackage{amsfonts}
\usepackage{nicefrac}
\usepackage{microtype}
\usepackage{bbm}

\usepackage{amsmath, amsfonts, amssymb}
\newtheorem{assumption}{Assumption}
\usepackage{mathtools}
\usepackage{graphicx}
\usepackage{bm}
\usepackage{lmodern}

\usepackage{physics}

\usepackage{algorithm}
\usepackage{algpseudocode}

% Common shortcuts
\def\mbb#1{\mathbb{#1}}
\def\mfk#1{\mathfrak{#1}}

\def\bN{\mbb{N}}
\def \C{\mbb{C}}
\def \R{\mbb{R}}
\def\bQ{\mbb{Q}}
\def\bZ{\mbb{Z}}
\def \cph{\varphi}
\renewcommand{\th}{\theta}
\def \ve{\varepsilon}
\newcommand{\mg}[1]{\left \| #1 \right\|}

% Often helpful macros
\newcommand{\floor}[1]{\left\lfloor#1\right\rfloor}
\newcommand{\ceil}[1]{\left\lceil#1\right\rceil}
\renewcommand{\P}{\mathbb P\qty}
\newcommand{\E}{\mathbb{E}\qty}
\newcommand{\Et}{\mathbb{E}_{\mathcal F_{t-1}} \qty}
\newcommand{\Pt}{\mathbb{P}_{\mathcal F_{t-1}} \qty}
\newcommand{\Cov}{\mathrm{Cov}\qty}
\newcommand{\Var}{\mathrm{Var}\qty}

\newcommand{\lr}{\mathrm{lr}}
\renewcommand{\grad}{\nabla} 
\newcommand{\1}{\mathbbm{1}}
\newcommand{\Mirr}{\mathrm{Mirr}}

% Sets
\usepackage{braket}

\graphicspath{{/}}
\usepackage{float}

\renewcommand{\set}[1]{\left\{#1\right\}}
\newcommand{\given}{\;\middle|\;}

\title{\bfseries Learning Rate for EGGROLL}
\date{\today}
\author{Rohan Mukherjee}

\begin{document}
\maketitle
\begin{assumption}
    $f \in \mathcal C^2(\R^d)$ is $L$-smooth,
    \begin{align*}
        \mg{\grad f(x) - \grad f(y)} \leq L \mg{x-y}
    \end{align*}
\end{assumption}

\section{Gradient Flow}
Recall the classic gradient descent update for an $L$-smooth function $f$
\begin{align}
    \label{eq:gradient-flow}
    x_{k+1} = x_k - \eta \grad f(x_k)
\end{align}
We want to turn this into an ODE. The natural time scale is $\Delta t = \eta$, which leads to the Ansatz $X(k \eta) = x_k$. Letting $t = k\eta$, we have $X_t = x_k$ and a taylor expansion shows $x_{k+1} - x_k = X(t + \eta) - X(t) \approx \dot X(t) \eta + o(\eta^2)$. This leads to 
\begin{align*}
    \dot X \eta + o(\eta^2) = -\eta \grad f(X)
\end{align*}
Matching terms with $\eta$ we get gradient flow
\begin{align*}
    \dot X_t + \grad f(X_t) = 0
\end{align*}
Now notice that 
\begin{align*}
    \dv{t} f(X_t) = \langle \grad f(X_t), \dot X_t \rangle = -\mg{\grad f(X_t)}^2
\end{align*}
So, 
\begin{align*}
    f(X_0) - f(X_T) = \int_0^t \mg{\grad f(X_t)}^2 \dd t
\end{align*}
While $|\grad f(X_t)| \geq \ve$, using $f(X_0) - f(X_T) \leq f(X_0) - f^*$, where $f^* = \inf_x f(x)$, this says
\begin{align*}
    T \ve^2 \leq f(X_0) - f^*
\end{align*}
So in $T = (f(X_0) -f^*) / \ve^2$ iterations we must have $\mg{\grad f(X_t)} < \ve$. This proof is so simple; what can learn from the continuous-time limit of EGGROLL?

\section{EGGROLL Gradient Flow}
Ignoring quadratic error, the (one-point) EGGROLL estimator becomes
\begin{align*}
    \tilde \grad f(x) = \frac{f(x + \sigma Z) - f(x - \sigma Z)}{2\sigma} Z = ZZ^\top \grad f(x)
\end{align*}
which is an unbiased estimator for $\grad f(X)$. Its covariance matrix is 
\begin{align*}
    \Cov(ZZ^\top \grad f(x)) = \mg{\grad f(x)}^2 I + \grad f(x) \grad f(x)^\top
\end{align*}
Let $\Sigma$ satisfy
\begin{align*}
    \Sigma \Sigma^\top = \mg{\grad f(x)}^2 I + \grad f(x) \grad f(x)^\top
\end{align*}
Writing $\tilde \grad f(x) = \grad f(x) + \xi$, we make the simplifying assumption that $\xi = \Sigma W$ for $W \sim \mathcal N(0,I)$ independent of $Z$. For EGGROLL descent
\begin{align*}
    x_{k+1} = x_k - \eta \tilde \grad f(x_k) = x_k - \eta (\grad f(x_k) + \xi_k)
\end{align*}
We can turn this into an SDE using the same Ansatz as above $X(t) = X(k \eta) = x_k$. To deal with the noise term, recall that for Brownian Motion $B_t$, we have $\Cov(\Sigma \Delta B_t) = \Sigma \Sigma^\top \Cov(\Delta B_t) = \Sigma \Sigma^\top \E[(\Delta B_t^1)^2]I = \eta \Sigma \Sigma^\top$. Also clearly $\Cov(\sqrt{\eta} \xi) = \eta \Sigma \Sigma^\top$. Thus we can write the SDE
\begin{align*}
    \dd X_t = -\grad f(X_t)\dd t + \sqrt{\eta} \Sigma \dd B_t
\end{align*}
Now notice by Ito's Lemma that
\begin{align*}
    \dd f(X_t) = \langle \grad f(X_t), \dd X_t \rangle + \frac \eta 2 \Tr(\grad^2 f(X_t) \Sigma \Sigma^\top) \dd t
\end{align*}
Therefore,
\begin{align*}
    \dd f(X_t) =& - \mg{\grad f(X_t)}^2 \dd t + \mathrm{martingale} 
    \\&+ \frac \eta 2 \qty( \mg{\grad f(X_t)}^2 \Tr(\grad^2 f(X_t)) + \grad f(X_t)^\top \grad^2 f(X_t) \grad f(X_t) )\dd t
\end{align*}
Under the $L$-smooth assumption, $\grad^2 f(x) \preceq LI$, so $\Tr(\grad^2 f) \leq Ld$, and $\grad f(X_t)^\top \grad^2 f(X_t) \grad f(X_t) \leq \mg{\grad f(X_t)}^2 L$. Therefore,
\begin{align*}
    df(X_t) = \qty(-1+\frac{\eta d(L+1)}{2}) \mg{\grad f(X_t)}^2 + \mathrm{martingale}
\end{align*}
Therefore, for $\E[f(X_t)]$ to decrease, we need $\eta < 2/d(L+1)$. This is very general, will work under the $L$-smoothness assumption, but it is more restrictive than necessary. We see that the necessary and sufficient condition on $\eta$ for function decrease is
\begin{align*}
    \boxed{\eta < \frac{2}{\Tr(\grad^2 f(X_t)) + \frac{\grad f(X_t)^\top \grad^2 f(X_t) \grad f(X_t)}{\mg{\grad f(X_t)}^2}}}
\end{align*}
For simplicity notice
\begin{align*}
    \frac{\grad f(X_t)^\top \grad^2 f(X_t) \grad f(X_t)}{\mg{\grad f(X_t)}^2} \leq \lambda_1(\grad^2 f(X_t)) \leq \Tr(\grad^2 f(X_t))
\end{align*}
Therefore a simpler necessary condition is
\begin{align*}
    \eta < \frac{1}{\Tr(\grad^2 f(X_t))}
\end{align*}
Of course, higher $\eta$ is equivalent to more function decrease when discretizing this SDE. So, if we could estimate $\Tr(\grad^2 f(X_t))$, we could use this as an adaptive learning rate, possibly speeding up convergence.


\section{Estimating the Trace of the Hessian}
Recall that the EGGROLL gradient is an unbiased estimator for the smoothed function $f_\mu(x) \coloneqq \E[f(x+\mu Z)]$,
\begin{align*}
    \grad f_\mu(x) = \frac{1}{\mu} \E[Zf(x+\mu Z)]
\end{align*}
The second-order Stein identity gives
\begin{align*}
    \grad^2 f_\mu(x) = \frac{1}{\mu^2} \E[(ZZ^\top - I) f(x+\mu Z)]
\end{align*}
Therefore,
\begin{align*}
    \Tr(\grad^2 f_\mu(x)) = \frac{1}{\mu^2} \E[(\mg{Z}^2 - d)f(x + \mu Z)]
\end{align*}
However, Taylor expanding, $f(x+\mu Z) = f(x) + \mu \langle \grad f(x), Z\rangle + O(\mu^2 \mg{Z}^2)$. It is clear that the bias coming from $1/\mu^2 f(x)$ and $1/\mu \langle \grad f(X), Z \rangle$ could totally destroy any useful signal this estimator will give us. 

However, there is a simple fix. Recall the simple fact that for any sufficiently nice function $g$
\begin{align*}
    \lim_{h \to 0} \frac{g(x+h) + g(x-h) - 2g(x)}{h^2} = g''(x)
\end{align*}
In particular, $g(x+h) + g(x-h) - 2g(x) = O(h^2)$. This is much better than the $\Omega(1)$ term $f(x+\mu Z)$ in the estimator. So, we propose using
\begin{align*}
    \Tr(\grad^2 f_\mu(x)) = \frac{1}{\mu^2} \E[(\mg{Z}^2 - d)(f(x + \mu Z) + f(x-\mu Z) - 2f(x))]
\end{align*}
\section{Addressing Batch Error}
In the simplest case of a neural network $\Phi_\theta: \R^d \to \R$ fitting the data $X \in \R^{N \times d}$ to $y \in \R^d$, we can use the MSE loss
\begin{align*}
    \ell_\theta(X) = \frac{1}{N} \sum_{i=1}^N \mg{\Phi_\theta(x_i) - y_i}^2
\end{align*}
Sampling a mini-batch $B \sim \binom{|N|}{k}$ we estimate $\ell$ as
\begin{align*}
    \ell_\theta(X ; B)=\frac{1}{|B|} \sum_{i \in B} \mg{\Phi_\theta(x_i) - y_i}^2
\end{align*}
More generally for an arbitrary loss function $\mathcal L$, we can write
\begin{align*}
    \ell_\theta(X) = \frac{1}{n} \sum_{i=1}^n \mathcal L(\Phi_\theta(x_i), y_i)
\end{align*}
and
\begin{align*}
    \ell_\theta(X;B) = \frac{1}{|B|} \sum_{i \in B} \mathcal L(\Phi_\theta(x_i), y_i)
\end{align*}
Write $\ell_\theta(X ; B) = \ell_\theta(X) + \xi_\theta(X ; B)$ where $\xi_\theta(X ; B)$ is the batch noise. We see immediately upon defining
\begin{align*}
    a_i = \begin{cases}
        \frac 1k - \frac 1n, \quad i \in B \\
        -\frac 1n, \quad i \not \in B
    \end{cases}
\end{align*}
that $\xi_\theta(X ; B) = \sum_{i=1}^n a_i \mathcal L(\Phi_\theta(x_i), y_i)$. 

Assuming $\mathcal L$ is $G$-Lipschitz and $\beta$-smooth,
\begin{align*}
    \mg{\grad \mathcal L(z, y)} \leq G, \qquad \mg{\grad \mathcal L(z,y) - \grad \mathcal L(z', y)} \leq \beta \mg{z-z'}
\end{align*} 
and assuming $\Phi_\theta$ is $K$-Lipschitz in $\theta$ and $L_\Phi$ smooth, 
\begin{align*}
    \mg{D_\theta \Phi_\theta(x)} \leq K, \qquad \mg{D_\theta \Phi_\theta(x) - D_\theta \Phi_\theta'(x)} \leq L_\Phi \mg{\theta - \theta'}
\end{align*}
each $\mathcal L(f_\theta(x_i), y_i)$ is $(GL_\Phi + \beta K^2)$-smooth in $\theta$, so the batch noise $\xi(X;B)$ is $2(GL_\Phi + \beta K^2)$-smooth. 

Recalling the descent inequality for $L$-smooth functions,
\begin{align*}
    f(x) + \langle \grad f(x), h \rangle - \frac{L}{2} \mg{h}^2 \leq f(x+h) \leq f(x) + \langle \grad f(x), h \rangle + \frac{L}{2} \mg{h}^2
\end{align*}
One can re-arrange this into the following central second difference bound
\begin{align*}
    |f(x+h) - f(x-h) - 2f(x)| \leq L |h|^2
\end{align*}
Therefore, the batch noise coming from
\begin{align*}
    \ell_{\theta + \sigma P}(X ; B) + \ell_{\theta - \sigma P}(X ; B) + 2 \ell_{\theta}(X ; B)
\end{align*}
is of order $O(\sigma^2)$. 
\section{Conclusion}
Using central differences is a powerful way to reduce variance in the EGGROLL estimator. Using second central differences can help us eliminate batch noise when trying to estimate the trace of the hessian. Estimating the trace of the hessian is important because it can be used as an adaptive, possibly better step size than just cross-validation would give. It could also eliminate the need for cross-validation for tuning the learning rate. We believe this could speed up training substantially. 
\end{document}